%% sections/06-conclusion.tex

\section{Conclusion}
\label{sec:conclusion}

NPC emotional moments in TRPG are often the sessions that players remember
longest --- and the hardest for adventure designers to produce reliably.
The conventional approach is to design rich NPC backstories and trust the
DM to deploy them at the right moment. This works when DMs are experienced,
attuned to player state, and operating with low cognitive load. It fails ---
silently, without obvious error --- when those conditions are not met.

The NPC arc-completion pattern addresses this failure mode by moving the
essential design work from runtime to design time. The question ``what should
this NPC say, and when?'' is answered before play begins: the arc-completion act
is specified verbatim, the firing condition is specified as a concrete observable
party action, and the DM discipline note specifies what must not happen after.
This pre-specification does not constrain play; it enables the party's actions
to become the cause of the NPC's most significant moment, which is a richer
design than any DM-authored monologue can produce.

The Marathon corpus documents 52 instances of this pattern firing correctly
across 49 sessions and 7 campaigns. Session coverage is 90\%. Zero instances
required DM scripting. The pattern holds for morally grey NPCs, professional-code
NPCs, information-holders, craft-witnesses, and situationally-triggered arcs
--- not only in the grief-adjacent campaigns where it was first documented.
The failure analysis confirms that NPCs without arc-completion conditions
produce memorable moments in fewer than 10\% of sessions where they are the
primary NPC encounter, even in high-scoring sessions.

Three design principles are directly actionable from this evidence:

\begin{enumerate}
  \item \textbf{Specify the act and the condition before play.}
    Every significant NPC should have an arc-completion act (the specific thing
    they are capable of saying or doing) and a firing condition (the specific
    observable party action that triggers it) written into their design document
    before the first session. If you cannot specify these two things, you have
    not finished designing the NPC.

  \item \textbf{Make it smaller than you think it should be.}
    The arc-completion text should be the shortest formulation that is semantically
    complete. If it is longer than two sentences, it is probably not the arc-completion;
    it is the designer's wish for what should follow. The accumulated context of
    prior sessions does the work; the arc-completion needs only to be the hinge.

  \item \textbf{Enforce the discipline note.}
    Do not deliver the arc-completion before the firing condition is met. Do not
    extend it afterward. Do not soften it for morally grey NPCs. The discipline
    note is not a suggestion; it is the mechanism by which the party's actions
    become the cause of the moment rather than the DM's timing.
\end{enumerate}

Future work should apply the pattern in human-run TRPG sessions and measure
DM discipline compliance, player-experience outcomes, and transfer to systems
with different NPC design conventions. The bilateral arc-completion sub-type
--- in which NPC and PC arc-completion conditions share a trigger domain and
fire simultaneously --- warrants its own study: the four bilateral instances
in this corpus produced the highest single-scene quality scores observed,
suggesting that the design of shared trigger domains between NPC and PC arcs
is a high-leverage design intervention that has not yet been systematically
investigated.

The claim this paper makes is modest and falsifiable: NPCs designed with
explicit arc-completion conditions produce memorable moments more reliably
than NPCs without them, and they do so without requiring the DM to script,
time, or improvise the moment. The corpus evidence supports the claim. The
design template is available, the discipline is learnable, and the failure
mode --- DM-scripted moments in place of character-logic moments --- is
identifiable and correctable. There is no excuse for NPCs that do nothing.
